\section*{IV.\quad ANALYSIS}
% ══════════════════════════════════════════════

{\color{rule}\hrule height 0.5pt}\vspace{6pt}

\subsection{Applicable Legal Standard}

\creac{Conclusion.}
Delaware courts apply the \textit{business judgment rule} as the default standard of review for board decisions, including merger approvals.  \textit{Aronson v.\ Lewis}, 473 A.2d 805, 812 (Del.\ 1984).  However, where a majority of the approving directors are interested or lack independence, or where the transaction was not the product of arm's-length bargaining, the more exacting \textit{entire fairness} standard applies.  \textit{Weinberger v.\ UOP, Inc.}, 457 A.2d 701, 710--11 (Del.\ 1983).

\creac{Rule.}
Under the entire fairness standard, the Board must demonstrate both (1) \textit{fair dealing}---encompassing the timing, initiation, structure, and approval of the transaction---and (2) \textit{fair price}---reflecting the economic and financial considerations of the transaction.  \textit{Id.}\ at 711.  The burden of proof rests with the defendant-directors once an interested transaction is shown.  \textit{Kahn v.\ Lynch Commc'n Sys., Inc.}, 638 A.2d 1110, 1117 (Del.\ 1994).

\creac{Application.}
Two of the nine directors---CEO Vance and CFO Song---have a material financial interest in the transaction not shared by common stockholders.  Under \textit{Rales v.\ Blasband}, 634 A.2d 927, 936 (Del.\ 1993), a director is ``interested'' when she will receive a personal financial benefit from a transaction that is not equally shared by the stockholders.  Because these two directors led the negotiations and will present the recommendation to the Board, there is a substantial likelihood that a court would apply entire fairness review unless the Board takes curative measures.

\creac{Conclusion.}
Absent curative measures, the business judgment rule's presumption is unavailable.  Entire fairness review will likely govern, materially increasing litigation exposure.

\hairrule

\subsection{The Absence of a Fairness Opinion}

\creac{Conclusion.}
While Delaware law does not categorically require a fairness opinion, the absence of one in a self-interested, compressed-timeline transaction of this size is a significant governance deficiency that courts and plaintiffs' counsel will scrutinize.

\creac{Rule.}
In \textit{In re Dollar Thrifty Shareholder Litigation}, 14 A.3d 573, 598 (Del.\ Ch.\ 2010), the court noted that a board need not obtain a fairness opinion as a matter of law but that its absence, combined with other process deficiencies, undermines the board's ability to demonstrate fair dealing.  Section 141(e) of the Delaware General Corporation Law permits directors to rely in good faith on information furnished by financial advisors, but this safe harbour presupposes the engagement of such an advisor.

\creac{Application.}
Here, four compounding factors heighten risk: (i) the \$2.3 billion transaction magnitude; (ii) director conflicts; (iii) absence of any market check or auction; and (iv) the accelerated timeline precluding customary due diligence.  In \textit{Lyondell Chemical Co.\ v.\ Ryan}, 970 A.2d 235, 243--44 (Del.\ 2009), the court held that Revlon duties required boards to act reasonably, not merely in good faith.  Failing to retain financial advisors when conflicts exist has been cited as evidence of unreasonable process.

\creac{Conclusion.}
We conclude that proceeding without a fairness opinion carries unacceptable risk of a successful entire fairness challenge and potential personal liability for directors.

\hairrule

\subsection{Revlon Duties and the Absence of a Market Check}

\creac{Conclusion.}
Because Greenfield stockholders will receive cash and surrender their equity interest, \textit{Revlon} duties are triggered and require the Board to seek the best price reasonably available.

\creac{Rule.}
\textit{Revlon, Inc.\ v.\ MacAndrews \& Forbes Holdings, Inc.}, 506 A.2d 173, 184 (Del.\ 1986), holds that when a company enters Revlon-mode (i.e., a sale of the entire company for cash), the board's role shifts from preserving the long-run entity to maximising short-term stockholder value.  While a formal auction is not mandatory, \textit{Barkan v.\ Amsted Indus., Inc.}, 567 A.2d 1279, 1286--87 (Del.\ 1989), the board must at least have reliable evidence that the price offered is the best available.

\creac{Application.}
No market check---formal or informal---was conducted.  The 28\% premium, while within the range of comparables, does not by itself satisfy \textit{Revlon} if the process was flawed.  In \textit{In re Toys ``R'' Us, Inc.\ Shareholder Litigation}, 877 A.2d 975, 1000 (Del.\ Ch.\ 2005), the court emphasised that a weak process can make even a premium price suspect.

\creac{Conclusion.}
The Board should authorise a discrete post-signing ``go-shop'' period or, preferably, conduct a targeted pre-signing outreach to two or three strategic buyers to establish a market-check record.

\vspace{10pt}

% ══════════════════════════════════════════════
